\section{{\v C}ech cohomology of quasi-coherent sheaves}
\label{section-cech-quasi-coherent}

\noindent
Let $X$ be a scheme.
Let $U \subset X$ be an affine open.
Recall that a {\it standard open covering} of $U$ is a covering
of the form $\mathcal{U} : U = \bigcup_{i = 1}^n D(f_i)$
where $f_1, \ldots, f_n \in \Gamma(U, \mathcal{O}_X)$ generate
the unit ideal, see
Schemes, Definition \ref{schemes-definition-standard-covering}.

\begin{lemma}
\label{lemma-cech-cohomology-quasi-coherent-trivial}
Let $X$ be a scheme.
Let $\mathcal{F}$ be a quasi-coherent $\mathcal{O}_X$-module.
Let $\mathcal{U} : U = \bigcup_{i = 1}^n D(f_i)$ be a standard
open covering of an affine open of $X$.
Then $\check{H}^p(\mathcal{U}, \mathcal{F}) = 0$ for
all $p > 0$.
\end{lemma}

\begin{proof}
Write $U = \Spec(A)$ for some ring $A$.
In other words, $f_1, \ldots, f_n$ are elements of $A$
which generate the unit ideal of $A$.
Write $\mathcal{F}|_U = \widetilde{M}$ for some $A$-module $M$.
Clearly the {\v C}ech complex
$\check{\mathcal{C}}^\bullet(\mathcal{U}, \mathcal{F})$
is identified with the complex
$$
\prod\nolimits_{i_0} M_{f_{i_0}} \to
\prod\nolimits_{i_0i_1} M_{f_{i_0}f_{i_1}} \to
\prod\nolimits_{i_0i_1i_2} M_{f_{i_0}f_{i_1}f_{i_2}} \to
\ldots
$$
We are asked to show that the extended complex
\begin{equation}
\label{equation-extended}
0 \to
M \to
\prod\nolimits_{i_0} M_{f_{i_0}} \to
\prod\nolimits_{i_0i_1} M_{f_{i_0}f_{i_1}} \to
\prod\nolimits_{i_0i_1i_2} M_{f_{i_0}f_{i_1}f_{i_2}} \to
\ldots
\end{equation}
(whose truncation we have studied in
Algebra, Lemma \ref{algebra-lemma-cover-module}) is exact.
It suffices to show that (\ref{equation-extended})
is exact after localizing at a prime $\mathfrak p$, see
Algebra, Lemma \ref{algebra-lemma-characterize-zero-local}.
In fact we will show that the extended complex localized
at $\mathfrak p$ is homotopic to zero.

\medskip\noindent
There exists an index $i$ such that $f_i \not \in \mathfrak p$.
Choose and fix such an element $i_{\text{fix}}$. Note that
$M_{f_{i_{\text{fix}}}, \mathfrak p} = M_{\mathfrak p}$. Similarly
for a localization at a product $f_{i_0} \ldots f_{i_p}$ and $\mathfrak p$
we can drop any $f_{i_j}$ for which $i_j = i_{\text{fix}}$.
Let us define a homotopy
$$
h :
\prod\nolimits_{i_0 \ldots i_{p + 1}}
M_{f_{i_0} \ldots f_{i_{p + 1}}, \mathfrak p}
\longrightarrow
\prod\nolimits_{i_0 \ldots i_p}
M_{f_{i_0} \ldots f_{i_p}, \mathfrak p}
$$
by the rule
$$
h(s)_{i_0 \ldots i_p} = s_{i_{\text{fix}} i_0 \ldots i_p}
$$
(This is ``dual'' to the homotopy in the proof of
Cohomology, Lemma \ref{cohomology-lemma-homology-complex}.)
In other words, $h : \prod_{i_0} M_{f_{i_0}, \mathfrak p} \to M_\mathfrak p$
is projection onto the factor
$M_{f_{i_{\text{fix}}}, \mathfrak p} = M_{\mathfrak p}$ and in general
the map $h$ equal projection onto the factors
$M_{f_{i_{\text{fix}}} f_{i_1} \ldots f_{i_{p + 1}}, \mathfrak p}
= M_{f_{i_1} \ldots f_{i_{p + 1}}, \mathfrak p}$. We compute
\begin{align*}
(dh + hd)(s)_{i_0 \ldots i_p}
& =
\sum\nolimits_{j = 0}^p
(-1)^j
h(s)_{i_0 \ldots \hat i_j \ldots i_p}
+
d(s)_{i_{\text{fix}} i_0 \ldots i_p}\\
& =
\sum\nolimits_{j = 0}^p
(-1)^j
s_{i_{\text{fix}} i_0 \ldots \hat i_j \ldots i_p}
+
s_{i_0 \ldots i_p}
+
\sum\nolimits_{j = 0}^p
(-1)^{j + 1}
s_{i_{\text{fix}} i_0 \ldots \hat i_j \ldots i_p} \\
& =
s_{i_0 \ldots i_p}
\end{align*}
This proves the identity map is homotopic to zero as desired.
\end{proof}

\noindent
The following lemma says in particular that for any affine scheme
$X$ and any quasi-coherent sheaf $\mathcal{F}$ on $X$ we have
$$
H^p(X, \mathcal{F}) = 0
$$
for all $p > 0$.

\begin{lemma}
\label{lemma-quasi-coherent-affine-cohomology-zero}
\begin{slogan}
Serre vanishing: Higher cohomology vanishes on affine schemes
for quasi-coherent modules.
\end{slogan}
Let $X$ be a scheme.
Let $\mathcal{F}$ be a quasi-coherent $\mathcal{O}_X$-module.
For any affine open $U \subset X$ we have
$H^p(U, \mathcal{F}) = 0$ for all $p > 0$.
\end{lemma}

\begin{proof}
We are going to apply
Cohomology, Lemma \ref{cohomology-lemma-cech-vanish-basis}.
As our basis $\mathcal{B}$ for the topology of $X$ we are going to use
the affine opens of $X$.
As our set $\text{Cov}$ of open coverings we are going to use the standard
open coverings of affine opens of $X$.
Next we check that conditions (1), (2) and (3) of
Cohomology, Lemma \ref{cohomology-lemma-cech-vanish-basis}
hold. Note that the intersection of standard opens in an affine is
another standard open. Hence property (1) holds.
The coverings form a cofinal system of open coverings of any element
of $\mathcal{B}$, see
Schemes, Lemma \ref{schemes-lemma-standard-open}.
Hence (2) holds.
Finally, condition (3) of the lemma follows from
Lemma \ref{lemma-cech-cohomology-quasi-coherent-trivial}.
\end{proof}

\noindent
Here is a relative version of the vanishing of cohomology of quasi-coherent
sheaves on affines.

\begin{lemma}
\label{lemma-relative-affine-vanishing}
Let $f : X \to S$ be a morphism of schemes.
Let $\mathcal{F}$ be a quasi-coherent $\mathcal{O}_X$-module.
If $f$ is affine then $R^if_*\mathcal{F} = 0$ for all $i > 0$.
\end{lemma}

\begin{proof}
According to
Cohomology, Lemma \ref{cohomology-lemma-describe-higher-direct-images}
the sheaf
$R^if_*\mathcal{F}$ is the sheaf associated to the presheaf
$V \mapsto H^i(f^{-1}(V), \mathcal{F}|_{f^{-1}(V)})$.
By assumption, whenever $V$ is affine we have that $f^{-1}(V)$ is
affine, see Morphisms, Definition \ref{morphisms-definition-affine}.
By Lemma \ref{lemma-quasi-coherent-affine-cohomology-zero} we conclude that
$H^i(f^{-1}(V), \mathcal{F}|_{f^{-1}(V)}) = 0$
whenever $V$ is affine. Since $S$ has a basis consisting of affine
opens we win.
\end{proof}

\begin{lemma}
\label{lemma-relative-affine-cohomology}
Let $f : X \to S$ be an affine morphism of schemes.
Let $\mathcal{F}$ be a quasi-coherent $\mathcal{O}_X$-module.
Then $H^i(X, \mathcal{F}) = H^i(S, f_*\mathcal{F})$ for all $i \geq 0$.
\end{lemma}

\begin{proof}
Follows from Lemma \ref{lemma-relative-affine-vanishing}
and the Leray spectral sequence. See
Cohomology, Lemma \ref{cohomology-lemma-apply-Leray}.
\end{proof}

\noindent
The following two lemmas explain when {\v C}ech cohomology
can be used to compute cohomology of quasi-coherent modules.

\begin{lemma}
\label{lemma-affine-diagonal}
Let $X$ be a scheme. The following are equivalent
\begin{enumerate}
\item $X$ has affine diagonal $\Delta : X \to X \times X$,
\item for $U, V \subset X$ affine open, the intersection
$U \cap V$ is affine, and
\item there exists an open covering $\mathcal{U} : X = \bigcup_{i \in I} U_i$
such that $U_{i_0 \ldots i_p}$ is affine open for all $p \ge 0$ and all
$i_0, \ldots, i_p \in I$.
\end{enumerate}
In particular this holds if $X$ is separated.
\end{lemma}

\begin{proof}
Assume $X$ has affine diagonal. Let $U, V \subset X$ be affine opens.
Then $U \cap V = \Delta^{-1}(U \times V)$ is affine. Thus (2) holds.
It is immediate that (2) implies (3). Conversely, if there is a
covering of $X$ as in (3), then $X \times X = \bigcup U_i \times U_{i'}$
is an affine open covering, and we see that
$\Delta^{-1}(U_i \times U_{i'}) = U_i \cap U_{i'}$
is affine. Then $\Delta$ is an affine morphism by
Morphisms, Lemma \ref{morphisms-lemma-characterize-affine}.
The final assertion follows from Schemes, Lemma
\ref{schemes-lemma-characterize-separated}.
\end{proof}

\begin{lemma}
\label{lemma-cech-cohomology-quasi-coherent}
Let $X$ be a scheme.
Let $\mathcal{U} : X = \bigcup_{i \in I} U_i$ be an open covering such that
$U_{i_0 \ldots i_p}$ is affine open for all $p \ge 0$ and all
$i_0, \ldots, i_p \in I$.
In this case for any quasi-coherent sheaf $\mathcal{F}$ we have
$$
\check{H}^p(\mathcal{U}, \mathcal{F}) = H^p(X, \mathcal{F})
$$
as $\Gamma(X, \mathcal{O}_X)$-modules for all $p$.
\end{lemma}

\begin{proof}
In view of
Lemma \ref{lemma-quasi-coherent-affine-cohomology-zero}
this is a special case of
Cohomology, Lemma
\ref{cohomology-lemma-cech-spectral-sequence-application}.
\end{proof}







